\documentclass[10pt,twocolumn]{article}
\usepackage[margin=0.75in]{geometry}
\usepackage{amsmath,amssymb}
\usepackage{booktabs}
\usepackage[numbers,sort&compress]{natbib}
\usepackage{hyperref}

% DRAFT. Title, author, affiliation, email, and submission date are the only
% remaining placeholders (in [square brackets]). Every number in Results,
% the Reward comparison, and Conclusions is copied from results/*.csv,
% named in-text at first use (see paper/literature_claims.md for citation
% verification status).

\title{[Paper title]}
\author{[Author name(s)]\\[Affiliation]\\[Email]}
\date{[arXiv submission date]}

\begin{document}
\maketitle

\begin{abstract}
Market makers earn the quoted spread while bearing inventory risk, and classical models such as Avellaneda--Stoikov and Gu\'eant--Lehalle--Fernandez-Tapia (GLFT) give closed-form quotes only under strong assumptions. We train QR-DQN and IQN with a CVaR tail-risk objective, applied at action selection rather than inside the Bellman update, inside an ABIDES-Gym limit-order-book simulator, and compare 17 agent--representation variants (5 agents; handcrafted, CNN, autoencoder, and recurrent state representations) against GLFT on held-out data across three volatility regimes. QR-DQN beats GLFT significantly in every regime ($p<0.0001$), a result that survives a reward-sign ablation and a tail-risk computation bug we found and fixed in IQN's CVaR action selection, after which IQN also beats GLFT in the regime where it previously lost. No learned state representation matches handcrafted features, and recurrent integration never significantly helps. The CVaR $\alpha$-sweep traces a return-versus-tail-risk frontier in direction at its extremes, but not cleanly in between at this seed count.
\end{abstract}

\noindent\textbf{Keywords:} distributional reinforcement learning; market making; CVaR; limit order book; agent-based simulation

\section{Introduction}
Market makers earn the quoted spread by providing liquidity, and bear inventory risk while they do so. Classical models give closed-form quotes under strong assumptions: Avellaneda and Stoikov \citep{avellaneda2008} derive quotes under a driftless Brownian reference price, and Gu\'eant, Lehalle and Fernandez-Tapia \citep{gueant2012} extend the model with hard inventory bounds. Reinforcement learning can, in principle, learn quotes without these assumptions, but the RL market-making work we review optimises expected reward and reports single policies \citep{spooner2018,kumar2020}.

Mean return alone hides how differently policies behave in the tail: in our own normal-regime results (Table~\ref{tab:main}), maximum drawdown ranges from $-\$141$ for the best policy to $-\$2353$ for the worst, a gap the Sharpe ratio alone does not show. Market making's downside comes from being caught with inventory during an adverse move, which is exactly a tail event, not an average one. Distributional RL \citep{bellemare2017} models the full return distribution rather than its mean, which makes a tail-risk criterion such as CVaR a direct fit, and is the objective we use throughout.

\subsection{Related Work}
\textbf{Analytical market making.} Avellaneda and Stoikov \citep{avellaneda2008} give a reservation-price model whose optimal spread depends on risk aversion and time to close. Gu\'eant et al.\ \citep{gueant2012} solve the same problem with hard inventory bounds, reducing it to an ODE system. We use their exact solution as the strongest analytical benchmark.

\textbf{RL for market making.} Spooner et al.\ \citep{spooner2018} apply temporal-difference learning with eligibility traces to a simulated limit order book, and report that plain Q-learning and SARSA fail in their basic setup while traces give profitable policies. They also introduce an inventory-dampened reward. Ganesh et al.\ \citep{ganesh2019} show that an RL market maker learns to skew quotes against inventory and to exploit drift in a multi-agent dealer market. Patel \citep{patel2018} uses a hierarchical pair of agents on Bitcoin order-book data, and reports that the multi-agent approach underperforms its macro-agent alone. Gu\'eant and Manziuk \citep{gueant2019} scale RL market making to portfolios of bonds. Ga\v{s}perov et al.\ \citep{gasperov2021survey} survey the area.

\textbf{Adversarial and risk-sensitive RL.} Spooner and Savani \citep{spooner2020} train a market maker against an adversary, and find that a strategic adversary dominates fixed or random ones. They note that risk-sensitive methods can suffer numerical instability, especially with exponential utility. Distributional RL learns the full return distribution rather than its mean \citep{bellemare2017}; quantile regression \citep{dabney2018qr} and implicit quantile networks \citep{dabney2018iqn} make this practical. Lim and Malik \citep{lim2022} show that the common distributional action-selection rule for CVaR need not converge to a CVaR-optimal policy, even when that policy is stationary and Markov.

\textbf{Simulation.} ABIDES \citep{byrd2019} is a discrete-event market simulator with nanosecond timestamps and a single seeded random-number generator, so that runs are reproducible. ABIDES-Gym \citep{amrouni2021} wraps it in the Gym interface. Hawkes-process order flow \citep{bacry2015} and queue-reactive models \citep{huang2015} are common models of order flow, which we do not use (see Section~\ref{sec:limits}).

\subsection{Contributions}
\begin{itemize}
  \item A CVaR action-selection rule for QR-DQN and IQN, evaluated across a sweep of tail levels $\alpha$ to trace a return-versus-tail-risk frontier.
  \item A comparison of 17 agent--representation variants (SARSA, plus DQN, PPO, QR-DQN and IQN with handcrafted, CNN, autoencoder and recurrent state representations) on one protocol.
  \item A held-out evaluation on 15 episodes per policy, with paired tests against GLFT, Avellaneda--Stoikov and a fixed-spread baseline, across three volatility regimes.
  \item A test of whether learned quotes recover the inventory-skew structure of the analytical models.
  \item A sensitivity check of the asymmetric reward's sign, showing the central GLFT result is robust to it.
\end{itemize}

\section{Methods}
\label{sec:methods}

\subsection{Market and simulator}
We use ABIDES-Gym \citep{amrouni2021} with the background agents of ABIDES's \texttt{rmsc04} configuration, unmodified. An episode lasts $T=390$ steps, equal to a 6.5-hour trading day at one step per minute. The agent quotes one lot of 100 shares per order, and inventory is bounded by $|q_t|\le Q_{\max}=1000$ shares. Prices are in dollars with tick size $0.01$.

\subsection{MDP}
We model market making as a finite-horizon MDP $(\mathcal S,\mathcal A,R,P,\gamma)$ with $\gamma=0.99$. The state is an 18-dimensional handcrafted feature vector of inventory, spread, mid-price move, order-book imbalance over three levels, signed volume, realised volatility, RSI and the remaining time fraction $\tau_t=(T-t)/T$; the components are listed in \texttt{encoders/handcrafted.py}. An action chooses a bid offset and an ask offset from the same grid,
\begin{equation}
  (\delta^b_t,\delta^a_t)\in\{0,10,\dots,100\}^2\ \text{ticks},\qquad |\mathcal A|=121,
\end{equation}
which quotes $p^b_t=m_t-0.01\,\delta^b_t$ and $p^a_t=m_t+0.01\,\delta^a_t$, where $m_t$ is the mid price.

\subsection{Reward}
\label{sec:reward}
Let $n^b_t,n^a_t$ be the quantities filled at the bid and ask in step $t$, and let $q_{t-1}$ be the inventory before the step's fills. The half-spread captures and the mark-to-market inventory P\&L are
\begin{equation}
  \psi^a_t=n^a_t\,(p^a_t-m_t),\quad
  \psi^b_t=n^b_t\,(m_t-p^b_t),\quad
  \mathrm{inv}_t=q_{t-1}\,(m_t-m_{t-1}),
\end{equation}
and $\mathrm{PnL}_t=\psi^a_t+\psi^b_t+\mathrm{inv}_t$. Our asymmetric reward penalises only adverse inventory P\&L:
\begin{equation}
  r_t=\mathrm{PnL}_t-\eta\,\max(0,-\mathrm{inv}_t),\qquad \eta=0.5.
  \label{eq:reward}
\end{equation}
At episode end a terminal penalty $-\kappa|q_T|\sigma_T$ is added, with $\kappa=1$ and $\sigma_T$ the standard deviation of the episode's log mid-price returns up to that step. The reward above is our own formulation. Section~\ref{sec:reward_comp} compares it with the reversed-sign form $\mathrm{PnL}_t-\eta\max(0,\mathrm{inv}_t)$, which is the functional form of Spooner et al.\ \citep{spooner2018}, implemented in our environment.

\subsection{Distributional agents}
\textbf{QR-DQN.} The return distribution of each action is represented by $N=32$ quantile atoms $\theta_i(s,a)$ at midpoints $\hat\tau_i=(2i-1)/(2N)$ \citep{dabney2018qr}. The target uses Double DQN action selection, $a^*=\arg\max_a \bar\theta(s',a)$ with $\bar\theta$ the mean over atoms, and the loss is the quantile Huber loss
\begin{equation}
  \mathcal L=\sum_{i=1}^{N}\frac1N\sum_{j=1}^{N}
  \big|\hat\tau_i-\mathbb 1\{u_{ij}<0\}\big|\,\mathcal L_\kappa(u_{ij}),\qquad
  u_{ij}=r+\gamma\theta_j(s',a^*)-\theta_i(s,a),
\end{equation}
where $\mathcal L_\kappa$ is the Huber loss with threshold $\kappa=1$. Both networks use a dueling head and prioritised replay.

\textbf{IQN.} The quantile function is a function of a sampled level $\tau$ \citep{dabney2018iqn}:
\begin{equation}
  Z_\theta(s,a,\tau)=f_\theta\big(\phi(\tau)\odot\psi(s),a\big),\qquad
  \phi(\tau)=\mathrm{ReLU}\Big(\sum_{j=0}^{31}\cos(\pi j\tau)\,w_j+b\Big),
\end{equation}
trained on $\tau\sim U[0,1]$, with the same dueling head and replay as QR-DQN.

\subsection{CVaR action selection}
\label{sec:cvar}
Risk enters only when an action is chosen, not in the Bellman update:
\begin{equation}
  a_t=\arg\max_{a}\ \mathrm{CVaR}_\alpha\big(Z(s_t,a)\big),\qquad
  \mathrm{CVaR}_\alpha(Z)\approx\frac{1}{n_\alpha}\sum_{i=1}^{n_\alpha}Z_{(i)},
\end{equation}
where $Z_{(i)}$ are the sorted atoms (QR-DQN) or sampled quantile values (IQN), and $n_\alpha=\max(1,\lfloor\alpha N\rfloor)$. This is the standard distributional action-selection rule \citep{dabney2018qr}. Lim and Malik \citep{lim2022} show that this rule need not converge to a CVaR-optimal policy, even when that policy is stationary and Markov. We therefore treat our wrapper as a selection heuristic and measure its effect empirically rather than claiming CVaR optimality.

For IQN, the $K$ samples are drawn from $U[0,\alpha]$ and their mean is the CVaR estimate, so no further tail restriction is applied. An earlier implementation applied the bottom-$\alpha$ slice a second time, which restricted the effective tail to $\alpha^2$ rather than $\alpha$; every IQN CVaR result in this paper is from the corrected implementation, retrained and re-evaluated after the fix.

\subsection{Recurrent variant}
The recurrent variant adds a 128-unit LSTM over the last $T=30$ observations inside the agent \texttt{agents/recurrent\_base.py}. Replay samples 30-step windows that lie within a single episode, and the hidden state is zeroed at the start of each window, following the random-update scheme of Hausknecht and Stone \citep{hausknecht2015}. An earlier version allowed windows to cross episode boundaries; all results reported here use the corrected sampler.

\subsection{Encoders}
The main comparison uses the handcrafted encoder. The CNN and autoencoder encoders, and the recurrent variant, are compared as ablations in Section~\ref{sec:results}. The autoencoder is pretrained on order-book snapshots and frozen during RL training.

\subsection{Analytical baselines}
Avellaneda--Stoikov \citep{avellaneda2008} quotes around the reservation price $r_t=m_t-q_t\gamma\sigma^2\hat\tau$, with $\hat\tau$ the remaining time, using a full spread of $\delta^\ast=\gamma\sigma^2\hat\tau+\frac{2}{\gamma}\ln(1+\frac{\gamma}{\kappa})$. The volatility $\sigma$ is an online estimate. GLFT \citep{gueant2012} is the exact solution of the inventory-bounded problem, with inventory limits $[-Q,Q]$ and $Q=10$ lots on the solver's grid. The half-spreads are
\begin{equation}
  \delta^b(q)=\frac{1}{\kappa}\ln\frac{v_q}{v_{q+1}}+\frac{\xi}{2}+\frac{1}{\gamma}\ln\Big(1+\frac{\gamma}{\kappa}\Big),
\end{equation}
with $\delta^a(q)$ defined symmetrically, where $(v_q)$ solves the linear ODE system of Proposition~5 in \citep{gueant2012}. We use the exact solution only; the linearised approximations are not used. A fixed-spread baseline quotes constant offsets.

\subsection{Evaluation}
Each policy is run greedily on 15 held-out episodes with seeds $80042$--$80056$, disjoint from the training seeds and from the seeds used to select checkpoints. The checkpoint for each run is the one with the highest training-time evaluation Sharpe. For each episode we compute the step P\&L $\mathrm{pnl}_t=\mathrm{spread\_pnl}_t+q_{t-1}(m_t-m_{t-1})$, which does not depend on the training reward. The intraday Sharpe ratio is
\begin{equation}
  \mathrm{Sharpe}=\frac{\mathrm{mean}_t(\mathrm{pnl}_t)}{\mathrm{std}_t(\mathrm{pnl}_t)}\sqrt{T},
\end{equation}
and the tail measure $\mathrm{CVaR}_{0.10}$ is the mean of the worst 10\% of episode P\&Ls. We compare policies with paired $t$-tests and Wilcoxon signed-rank tests over the 15 seeds.

\section{Experiments and Results}
\label{sec:results}

\subsection{Protocol}
We evaluate five agents (SARSA, DQN, PPO, QR-DQN, IQN) against three analytical baselines (fixed-spread, Avellaneda--Stoikov, GLFT) in each of three regimes (\texttt{low\_vol}, \texttt{normal}, \texttt{high\_vol}), using the held-out protocol of Section~\ref{sec:methods} (15 unseen episodes, seeds $80042$--$80056$). QR-DQN and IQN use $\alpha=0.25$ unless swept. All reported $p$-values are from the paired Wilcoxon signed-rank test on the 15 shared seeds; the paired $t$-test agrees on every comparison below.

\subsection{Main results}
\label{sec:main-results}
Table~\ref{tab:main} reports the normal-regime results. QR-DQN has the highest held-out Sharpe of any policy, including the baselines, and the lowest downside tail risk ($\mathrm{CVaR}_{0.10}$) of the five RL agents.

\begin{table}[t]
\centering\footnotesize
\begin{tabular}{lrrrr}
\toprule
Model & Sharpe & Mean P\&L & CVaR$_{0.10}$ & MDD \\
\midrule
QR-DQN       & $+4.19$ & $+1298$ & $+617$  & $-141$ \\
DQN          & $+3.65$ & $+2146$ & $+1250$ & $-304$ \\
Fixed-spread & $+1.35$ & $+275$  & $-2340$ & $-652$ \\
SARSA        & $+0.11$ & $-95$   & $-2516$ & $-910$ \\
IQN          & $+0.01$ & $-86$   & $-2215$ & $-1419$ \\
PPO          & $-0.02$ & $-145$  & $-2237$ & $-1803$ \\
AS           & $-0.90$ & $-1185$ & $-3749$ & $-2199$ \\
GLFT         & $-1.17$ & $-1405$ & $-3618$ & $-2353$ \\
\bottomrule
\end{tabular}
\caption{Held-out performance, normal regime, 15 episodes (\texttt{results/holdout\_eval.csv}). CVaR$_{0.10}$ and MDD are means over episodes, in dollars.}
\label{tab:main}
\end{table}

Table~\ref{tab:glft} gives the paired significance test against GLFT in all three regimes. QR-DQN beats GLFT significantly in every regime ($p<0.0001$). IQN beats GLFT significantly in \texttt{normal} and \texttt{high\_vol}, and, after the CVaR fix of Section~\ref{sec:cvar}, also in \texttt{low\_vol} ($p=0.0001$) --- before the fix, IQN lost to GLFT in that regime. DQN, PPO, and SARSA beat GLFT in \texttt{normal} and \texttt{high\_vol} but not \texttt{low\_vol}.

\begin{table}[t]
\centering\footnotesize
\begin{tabular}{llrrl}
\toprule
Regime & Agent & Mean $\Delta$ & $p$ & Beats GLFT \\
\midrule
low\_vol  & QR-DQN & $+1.61$ & $<0.0001$ & yes \\
low\_vol  & IQN    & $+1.79$ & $0.0001$  & yes \\
low\_vol  & SARSA  & $-0.66$ & $0.095$   & no \\
low\_vol  & DQN    & $-1.20$ & $0.003$   & no \\
low\_vol  & PPO    & $-1.11$ & $0.0009$  & no \\
normal    & QR-DQN & $+5.36$ & $<0.0001$ & yes \\
normal    & DQN    & $+4.82$ & $<0.0001$ & yes \\
normal    & SARSA  & $+1.28$ & $<0.0001$ & yes \\
normal    & IQN    & $+1.17$ & $0.0009$  & yes \\
normal    & PPO    & $+1.15$ & $0.018$   & yes \\
high\_vol & DQN    & $+5.11$ & $<0.0001$ & yes \\
high\_vol & QR-DQN & $+3.97$ & $<0.0001$ & yes \\
high\_vol & IQN    & $+2.03$ & $0.003$   & yes \\
high\_vol & SARSA  & $+1.07$ & $0.026$   & yes \\
high\_vol & PPO    & $+1.54$ & $0.003$   & yes \\
\bottomrule
\end{tabular}
\caption{Mean Sharpe difference vs.\ GLFT, paired Wilcoxon $p$ (\texttt{results/significance\_vs\_glft.csv}).}
\label{tab:glft}
\end{table}

\subsection{CVaR efficient frontier}
We sweep $\alpha\in\{0.05,0.10,0.25,0.50,1.0\}$ for QR-DQN and IQN in the normal regime, with 4 seeds per (agent, $\alpha$) pair (\texttt{results/cvar\_efficient\_frontier.csv}). Only the two extremes of IQN's frontier behave as the method predicts: $\alpha=0.05$ has the lowest mean P\&L (\$4) and the least negative $\mathrm{CVaR}_{0.10}$ ($-\$1283$) of its five settings, and $\alpha=1.0$, the risk-neutral limit, has the most negative $\mathrm{CVaR}_{0.10}$ ($-\$2731$). The interior values are not monotonic: $\alpha=0.10$ has both a higher mean P\&L (\$413) and a less negative $\mathrm{CVaR}_{0.10}$ ($-\$696$) than $\alpha=0.25$ or $\alpha=0.50$. QR-DQN's frontier is non-monotonic throughout, including at the endpoints: $\alpha=0.25$ has both the highest mean P\&L (\$1000) and a positive $\mathrm{CVaR}_{0.10}$ ($+\$445$), better on both axes than the nominally more conservative $\alpha=0.05$ (\$639, $-\$880$). Per-seed standard deviations (4 seeds) are large relative to the differences between adjacent $\alpha$ values. We read this as a frontier that exists in direction at the extremes but is not resolved at this seed count, not as a null result.

\subsection{Out-of-distribution transfer}
Each agent's checkpoint, trained in one regime, is rolled out with no fine-tuning in the other two regimes, across all 6 directed regime pairs (\texttt{results/ood\_transfer\_summary.csv}). QR-DQN stays profitable in all 6 transfers and improves on its in-distribution Sharpe in 3 of them. DQN stays profitable in 5 of 6. IQN stays profitable in 4 of 6 and improves in 3. PPO and SARSA are the least robust, profitable in only 2 and 3 of 6 respectively. We do not report percentage degradation, since several in-distribution Sharpe values are close to zero and the ratio is not meaningful there (Section~\ref{sec:limits}).

\subsection{State representation}
\label{sec:encoder-results}
Replacing the handcrafted encoder with a CNN or a pretrained autoencoder over raw LOB depth, for all four neural agents, gives a clean negative result for two of them: handcrafted significantly outperforms both learned encoders for QR-DQN and DQN ($p<0.001$ in all four comparisons; QR-DQN: $+4.19$ vs.\ $+0.42$ and $+0.26$; DQN: $+3.65$ vs.\ $+0.20$ and $-0.63$), and shows no significant difference for PPO or IQN, whose own handcrafted baselines are themselves weak (\texttt{results/encoder\_ablation\_significance.csv}). This comparison only holds once the CNN and autoencoder encoders are given the same six private features (inventory, quote distances, outstanding offsets, time remaining) the handcrafted encoder always had; an earlier implementation omitted them, which we treat as a bug rather than a design choice, since the task's inventory-skew mechanism depends on the agent observing its own inventory.

The recurrent variant --- an LSTM(128) in place of the snapshot encoder, same features --- never significantly improves any of the ten agent/regime/$\alpha$ combinations we test, and significantly hurts QR-DQN, DQN, and IQN in the normal regime (\texttt{results/recurrent\_variant\_significance.csv}). Its interaction with the CVaR tail level is not consistent: recurrent QR-DQN is significantly worse at $\alpha=0.10$ but significantly better at $\alpha=0.50$.

\subsection{Validation}
AS-recovery (quote-skew $R^2$ against the GLFT closed form) is mixed: IQN shows strong recovery in \texttt{low\_vol} ($R^2=0.80$) and \texttt{high\_vol} ($R^2=0.70$), QR-DQN and DQN show weak recovery ($R^2\approx0.3$--$0.5$) in most regimes, and PPO never fills enough episodes to compute the statistic (\texttt{results/as\_recovery.csv}). The simulator passes 14 of 15 stylized-facts checks (fat tails, volatility clustering, price impact, queue-imbalance predictability, spread autocorrelation, across three regimes); only spread autocorrelation fails, in \texttt{low\_vol} (\texttt{results/stylized\_facts.json}).

\subsection{Reward comparison}
\label{sec:reward_comp}
We compare our reward (Section~\ref{sec:reward}) with its reversed-sign form, the functional form of Spooner et al.\ \citep{spooner2018}, across all 5 agents and 3 regimes at seed 42, plus 3 extra seeds for QR-DQN and IQN in \texttt{normal} (\texttt{results/reward\_comparison\_summary.csv}). Across the 15 agent$\times$regime pairs, ours is significantly better in 4, the reversed sign is significantly better in 4, and the remaining 7 show no significant difference --- neither sign dominates. The result that matters most for this paper is robust to the choice: QR-DQN beats GLFT significantly in the normal regime under both ours ($+5.36$, $p<0.0001$) and the reversed sign ($+3.73$, $p=0.0001$). The three extra seeds show the pairwise comparison itself can flip: IQN in \texttt{normal} favours ours significantly at two of three extra seeds and is tied at the third, while QR-DQN's extra seeds show no significant difference at any of the three.

\section{Conclusions and limitations}
\label{sec:limits}

Three findings. First, QR-DQN with CVaR action selection beats the exact GLFT solution on held-out data in every regime we test ($p<0.0001$), and the margin survives both the choice of reward sign (Section~\ref{sec:reward_comp}) and, for IQN, a tail-risk computation bug we found and fixed (Section~\ref{sec:cvar}): before the fix IQN lost to GLFT in \texttt{low\_vol}; after it, IQN also beats GLFT there. Second, no state representation we tried matches handcrafted features: CNN and autoencoder encoders over raw LOB depth lose significantly to handcrafted for QR-DQN and DQN, and an LSTM backbone never significantly helps any agent and significantly hurts three of them. Third, the CVaR $\alpha$-sweep produces a frontier that exists in direction at its extremes but is not cleanly monotonic at this training budget and seed count.

Limitations. The simulator's background order flow is ABIDES's own \texttt{rmsc04} configuration, unmodified; it is not calibrated to LOBSTER or any other real exchange data, and the stylized-facts check in Section~\ref{sec:main-results} validates internal realism, not a fit to a specific market. We do not model Hawkes-process (self-exciting) order arrivals \citep{bacry2015}, and we do not train against an adversary \citep{spooner2020}. Our CVaR wrapper is a selection heuristic, not a provably CVaR-optimal policy (Section~\ref{sec:cvar}); Lim and Malik's result \citep{lim2022} applies directly to it. We trade a single asset. Most ablations use one seed; the CVaR sweep and the reward comparison's QR-DQN/IQN rows use four, and even there the reward-sign comparison shows the pairwise result can flip between seeds. Percentage OOD degradation is not reported because several in-distribution Sharpe values are close to zero.

Future work: calibrating the background agents to LOBSTER data, Hawkes-process arrivals, adversarial training, and a CVaR objective applied inside the Bellman update rather than at selection, following a time-consistent formulation such as Lim and Malik's \citep{lim2022}.

\bibliographystyle{plainnat}
\bibliography{references}

\appendix
\section{Appendix}

\subsection{Full 17-variant grid}
\label{sec:appendix-grid}
Table~\ref{tab:grid} gives held-out Sharpe, normal regime, for every (agent, state representation) combination in the ablation grid: SARSA (handcrafted features only --- it is a from-scratch tile-coding agent, incompatible with the neural encoders of Section~\ref{sec:methods}) and the four neural agents, each with the handcrafted encoder, the CNN encoder, the autoencoder encoder, and the recurrent (LSTM) variant in place of a snapshot encoder. These are the same numbers reported piecemeal in Sections~\ref{sec:main-results} and \ref{sec:encoder-results}, collected here in one table.

\begin{table}[t]
\centering\footnotesize
\begin{tabular}{lrrrr}
\toprule
Agent & Handcrafted & CNN & Autoencoder & Recurrent \\
\midrule
SARSA   & $+0.11$ & --- & --- & --- \\
DQN     & $+3.65$ & $+0.20$ & $-0.63$ & $+0.35$ \\
PPO     & $-0.02$ & $-0.50$ & $-0.11$ & $-0.05$ \\
QR-DQN  & $+4.19$ & $+0.42$ & $+0.25$ & $+1.34$ \\
IQN     & $+0.01$ & $-0.05$ & $-0.14$ & $-0.63$ \\
\bottomrule
\end{tabular}
\caption{Held-out Sharpe, normal regime, all 17 variants (\texttt{results/holdout\_summary.csv}, \texttt{results/encoder\_ablation\_summary.csv}, \texttt{results/recurrent\_variant\_significance.csv}). SARSA has no CNN/autoencoder/recurrent cell: it is a tile-coding agent with no neural encoder to swap (Section~\ref{sec:methods}). QR-DQN's recurrent cell uses $\alpha=0.25$, matching every other cell in this table.}
\label{tab:grid}
\end{table}

\subsection{Hyperparameters}
\label{sec:appendix-hparams}
All values are from \texttt{training/configs/} and apply to every reported run unless noted otherwise.

\begin{table}[t]
\centering\footnotesize
\begin{tabular}{lll}
\toprule
Component & Parameter & Value \\
\midrule
All agents    & $\gamma$                     & $0.99$ \\
              & optimizer                     & Adam, lr $=10^{-4}$ \\
              & batch size                    & $256$ \\
              & replay buffer                 & $10^5$, prioritized \\
              & target update                 & every $1000$ steps \\
              & $\epsilon$-greedy             & $1.0\to0.05$ over $5\times10^4$ steps \\
QR-DQN        & quantile atoms $N$             & $32$ \\
              & hidden dim                    & $256$, dueling head \\
              & quantile-Huber $\kappa$        & $100$ \\
IQN           & quantile samples $K$           & $16$ \\
              & cosine embedding dim           & $32$ \\
              & hidden dim                     & $256$, dueling head \\
DQN           & hidden dim                     & $256$ \\
              & Huber $\beta$                  & $100$ \\
PPO           & hidden dim                     & $128$ \\
              & GAE $\lambda$                  & $0.95$ \\
              & clip target $\delta$           & $0.01$, $4$ epochs/update \\
              & entropy coef.\                 & $0.01$ \\
              & value coef.\                   & $0.005$ \\
SARSA         & tilings / tiles                & $16 / 8$ \\
              & $\alpha,\gamma,\lambda$        & $0.001, 0.97, 0.96$ \\
CVaR policy   & $\alpha$ (tail fraction)        & $0.25$ (sweep: $0.05,0.10,0.25,0.50,1.0$) \\
CNN encoder   & conv channels                  & $32\to16$, kernel $3$ \\
              & latent dim                     & $8$ \\
Autoencoder   & latent dim                     & $8$ (capacity-matched to CNN) \\
              & pretraining                    & frozen at RL-train time \\
Recurrent     & LSTM hidden / layers            & $128 / 1$ \\
              & window $T$                     & $30$ steps \\
Training      & episode budget                 & $2000$ episodes \\
              & episode length                 & $390$ steps ($6.5$h, 1 step/min) \\
              & eval cadence                   & every $50$ episodes, $10$ episodes \\
Environment   & action grid                    & $11\times11$ tick offsets $\{0,10,\dots,100\}$\textcent \\
              & $Q_{\max}$ / order size         & $1000 / 100$ shares \\
              & tick size                      & \$0.01 \\
\bottomrule
\end{tabular}
\caption{Hyperparameters, from \texttt{training/configs/\{agent,encoder,variant,env,config\}/*.yaml}.}
\label{tab:hparams}
\end{table}

\subsection{Reward formulations}
\label{sec:appendix-rewards}
All reported results use the asymmetric reward of Section~\ref{sec:reward} (ours, primary) and, in Section~\ref{sec:reward_comp} only, its reversed-sign comparator. Two further formulations are implemented in \texttt{envs/lob\_env.py} (\texttt{training/configs/reward/\{quadratic,sparse\}.yaml}) but were not run in this campaign; we give them here for completeness, since the code treats \texttt{reward\_type} as a first-class config axis:
\begin{itemize}
\item \textbf{Quadratic} (Cartea--Jaimungal flavour): $r_t = \mathrm{PnL}_t - \lambda q_{t-1}^2$, a symmetric penalty on inventory level rather than on adverse inventory P\&L, with $\lambda=0.001$.
\item \textbf{Sparse}: $r_t=+1$ on a round-trip fill (both sides filled), $r_t=-0.5$ on an adverse cross-spread fill, $r_t=0$ otherwise --- a diagnostic reward with no continuous P\&L signal.
\end{itemize}

\subsection{Full held-out results, all regimes}
\label{sec:appendix-holdout}
Table~\ref{tab:main} in Section~\ref{sec:main-results} reports the normal regime only. Tables~\ref{tab:lowvol} and \ref{tab:highvol} give the same metrics for \texttt{low\_vol} and \texttt{high\_vol} (\texttt{results/holdout\_eval.csv}, \texttt{results/holdout\_summary.csv}).

\begin{table}[t]
\centering\footnotesize
\begin{tabular}{lrrrr}
\toprule
Model & Sharpe & Mean P\&L & CVaR$_{0.10}$ & MDD \\
\midrule
IQN          & $+3.03$ & $+423$  & $+124$  & $-65$   \\
QR-DQN       & $+2.85$ & $+955$  & $+611$  & $-186$  \\
AS           & $+1.24$ & $+540$  & $+78$   & $-577$  \\
GLFT         & $+1.24$ & $+489$  & $-222$  & $-521$  \\
Fixed-spread & $+0.90$ & $+129$  & $-23$   & $-137$  \\
SARSA        & $+0.58$ & $+101$  & $-123$  & $-307$  \\
PPO          & $+0.13$ & $+104$  & $-258$  & $-755$  \\
DQN          & $+0.04$ & $+67$   & $-1131$ & $-1350$ \\
\bottomrule
\end{tabular}
\caption{Held-out performance, \texttt{low\_vol} regime, 15 episodes (\texttt{results/holdout\_eval.csv}). CVaR$_{0.10}$ and MDD are means over episodes, in dollars.}
\label{tab:lowvol}
\end{table}

\begin{table}[t]
\centering\footnotesize
\begin{tabular}{lrrrr}
\toprule
Model & Sharpe & Mean P\&L & CVaR$_{0.10}$ & MDD \\
\midrule
DQN          & $+2.93$ & $+5027$ & $+30$    & $-1532$ \\
QR-DQN       & $+1.79$ & $+2359$ & $-475$   & $-1604$ \\
IQN          & $-0.15$ & $-581$  & $-5834$  & $-3678$ \\
PPO          & $-0.64$ & $-1628$ & $-8808$  & $-4566$ \\
SARSA        & $-1.10$ & $-2563$ & $-7188$  & $-4460$ \\
AS           & $-2.18$ & $-4902$ & $-11429$ & $-5806$ \\
Fixed-spread & $-2.18$ & $-4902$ & $-11429$ & $-5806$ \\
GLFT         & $-2.18$ & $-4902$ & $-11429$ & $-5806$ \\
\bottomrule
\end{tabular}
\caption{Held-out performance, \texttt{high\_vol} regime, 15 episodes (\texttt{results/holdout\_eval.csv}).}
\label{tab:highvol}
\end{table}

Read together with Table~\ref{tab:main}, the three regimes show the baselines' Sharpe ranking is not stable across volatility --- AS and GLFT are the best performers in \texttt{low\_vol} and the worst in \texttt{normal} and \texttt{high\_vol} --- while QR-DQN is the only agent to finish in the top two of every regime.

\end{document}
